\section{Cancellation from square-difference-freeness: statement of the kernel estimate}\label{sec:kernel}

\paragraph{What this section proves and how to read it.}
Square-difference-freeness of $A$ produces cancellation in exactly one place
of the analytic stage: a bound for a truncated Fourier pairing of two functions
that live on single residue classes of a modulus $D$, on ordered intervals,
with no difference a positive square. This section states that bound;
elsewhere, in the rescaling of \cref{sec:scale-comparison}, only the fact that
the property passes to progressions with square step is used. The bound is
applied in
the accepted-pair cancellation (\cref{lem:record:density-can}), and through it
in the wrap of the signed index form (\cref{lem:record:wrap-expansion}). The
proof, a quadratic circle-method argument, is in \cref{app:kernel20}. The
statement carries symbolic exponents because the cutoff of the functional and
the range of the kernel have to match: the cancellation applies the kernel at
the cutoff $H=N^{1/40}$ with moduli up to $N^{1/20}$, and a kernel bound proved
for a smaller cutoff does not apply to it. The proof is the same for every
admissible choice of exponents, so it is given once.

\begin{definition}[Kernel notation]\label{def:kernel-notation}
For $\xi\in\Q/\Z$ let $q(\xi)$ be its reduced denominator, and let
\[
 g(\xi)=\frac1{\varphi(q(\xi))}\sum_{m\in(\Z/q(\xi)\Z)^\times}e\big(\xi m^2\big),
 \qquad g(0)=1,
\]
be the \emph{unit quadratic average} at $\xi$; thus $|g(\xi)|\le1$, and for
squarefree odd $q(\xi)$ this is the product formula of
\cref{lem:square-shift}(d). For $H\ge1$ and any positive integer $D$ put
\[
 \Xi_{H,D}=\{\xi\in\Q/\Z:\ q(\xi)\le H,\ q(\xi)\ \text{squarefree},\ (q(\xi),D)=1\}.
\]
When $D=8d$ with $d$ odd, strict square classes modulo $D$ are those of
\cref{def:strict-square-class};
by the paragraph after that definition they are exactly the squares of units
modulo $D$, so a strict square class is the square of a unit modulo every odd
prime power dividing $D$. For an integer $N\ge1$
and a finitely supported $F\colon\Z\to\C$ write
$\widehat F(\xi)=N^{-1}\sum_nF(n)e(-n\xi)$ and
$\|F\|_{\ell^2}=(\sum_n|F(n)|^2)^{1/2}$.
\end{definition}

Three exponents enter the proof: the modulus satisfies $D\lesssim N^\kappa$,
the cutoff is $H\asymp N^{\beta/2}$, and the major arcs of the circle method
are the neighbourhoods (of the radius fixed in \cref{app:kernel20}) of the
fractions with denominator at most $N^\eta$; the
saving is $H^{-1/3}\asymp N^{-\beta/6}$.

\begin{definition}[Kernel exponents]\label{def:kernel-conditions}
A triple $(\beta,\kappa,\eta)$ of positive numbers is \emph{admissible for the
kernel} if the following seven strict inequalities hold:
\begin{align*}
 &\text{(E1)}\quad \kappa+\tfrac\beta2-\tfrac\eta2<-\tfrac\beta6, &
 &\text{(E2)}\quad \kappa+\beta-\eta<-\tfrac\beta6, &
 &\text{(E3)}\quad \eta+2\kappa+\tfrac{3\beta}2-1<-\tfrac\beta6,\\
 &\text{(E4)}\quad 2\eta+2\kappa+\beta-\tfrac12<-\tfrac\beta6, &
 &\text{(K0)}\quad 3\eta<1, &
 &\text{(K1)}\quad \kappa+\tfrac\beta2<\eta,\\
 &\text{(K2)}\quad 2\eta+\kappa+\tfrac\beta2<1 .
\end{align*}
\end{definition}

Here (K0) makes the major arcs disjoint, (K1) puts every centre of the model
kernel (\cref{lem:record:K20-centres}) on a major arc, and (K2) separates
distinct centres; (E1)--(E4) say that the
four error terms of the circle method (the actual kernel on the minor arcs, the
model kernel on the minor arcs, and the model and the actual kernel on the
major arcs) are smaller than the saving $N^{-\beta/6}$.

\begin{proposition}[The actual-square kernel]\label{prop:record:kernel20}
Let $(\beta,\kappa,\eta)$ be admissible for the kernel, and fix
$C_D,C_H\ge1$. There are $C$ and $N_0$, depending only on $\beta$, $\kappa$,
$\eta$, $C_D$ and $C_H$, with the following property. Let $N\ge N_0$, let
$D=8d\le C_DN^\kappa$ with $d$ odd (not necessarily squarefree), and let
$C_H^{-1}N^{\beta/2}\le H\le C_HN^{\beta/2}$. Let $F,J\colon\Z\to\C$ be supported
in $[1,N]$, in residue classes $a$ and $b$ modulo $D$. Suppose that $x<y$ for
all $x\in\supp F$ and $y\in\supp J$, and that no such difference $y-x$ is the
square of a positive integer.
\begin{enumerate}[label=(\roman*)]
\item If $b-a$ is a strict square class modulo $D$, then
 \begin{equation}\label{eq:record:K20}
  \Big|D\sum_{\xi\in\Xi_{H,D}}g(\xi)\widehat F(-\xi)\widehat J(\xi)\Big|
  \le C\,N^{-\beta/6}\,\frac{\|F\|_{\ell^2}\|J\|_{\ell^2}}N .
 \end{equation}
\item Let $\mathcal P_0$ be a finite set of odd primes and $m_p\ge1$
 $(p\in\mathcal P_0)$ integers such that the exponent of $p$ in $D$ is
 exactly $2m_p$. Suppose that
 $b-a\equiv1\pmod8$, that $b-a$ is a nonzero square modulo every prime of $d$
 outside $\mathcal P_0$, and that at each $p\in\mathcal P_0$ either $b-a$ is a
 unit square modulo $p$, or $v_p(b-a)=2j_p$ with $1\le j_p<m_p$ and
 the unit part of the representative a nonzero square modulo $p$ (where $v_p(b-a)$ is
 the valuation of any integer representative of the class $b-a$
 modulo $p^{2m_p}$, independent of the choice when it is
 $<2m_p$). Then \eqref{eq:record:K20} holds
 with the same $C$ and $N_0$ as in (i), as a special case of (iii).
 The corresponding rescaling identity is as follows. With $R=\prod_pp^{2j_p}$
 over the primes of $\mathcal P_0$ at which $v_p(b-a)\ge2$, $D_R=D/R$ and
 $N'=\lceil N/R\rceil$, the expression inside the absolute value
 in \eqref{eq:record:K20} (the \emph{pairing}) equals
 $R(N'/N)^2$ times the corresponding expression with $D$ replaced
 by $D_R$, at length $N'$, for the functions $F',J'$ on $[1,N']$
 constructed in \cref{lem:record:nonunit-rescale}, to which (i)
 applies with the constants $C_DR^\kappa$ and $C_HR^{\beta/2}$.
\item With the condition $D=8d$, $d$ odd, replaced by
 $D=2^{h_2}\prod_{p\text{ odd}}p^{h_p}\le C_DN^\kappa$, $h_2\ge3$,
 \eqref{eq:record:K20} holds whenever $b-a$ is a nonvanishing square class
 (\cref{def:nonvanishing-square-class}). The constants $C$ and
 $N_0$ may be taken depending only on $\beta,\kappa,\eta,C_D,C_H$,
 as in (i). The root-class construction of
 \cref{lem:kernel-depth-classes} supplies for these classes the same
 arithmetic coefficient and the same period and Fourier-coefficient bounds
 as for the strict classes of (i), for the proof in \cref{app:kernel20}.
\end{enumerate}
In particular, if $|F|,|J|\le1$ and $F,J$ are supported in intervals
$I_1,I_2$, the right-hand side of \eqref{eq:record:K20} is at most
$CN^{-\beta/6}N^{-1}((|I_1|/D+1)(|I_2|/D+1))^{1/2}$.
\end{proposition}

The paper applies the proposition at one triple.

\begin{corollary}[The kernel at $\beta=1/20$]\label{cor:kernel20}
The triple $(\beta,\kappa,\eta)=(1/20,1/20,17/100)$ is admissible for the
kernel. The slacks of (E1), (E2), (E3), (E4), (K0), (K1), (K2), that is, the
differences between the right-hand and the left-hand sides, are
\[
 \frac1{600},\qquad\frac{37}{600},\qquad\frac{97}{150},\qquad\frac1{600},
 \qquad\frac{49}{100},\qquad\frac{19}{200},\qquad\frac{117}{200}.
\]
Hence \cref{prop:record:kernel20} holds with $D\le C_DN^{1/20}$,
$C_H^{-1}N^{1/40}\le H\le C_HN^{1/40}$ and the saving $N^{-\beta/6}=N^{-1/120}$.
These seven numbers are (S6) of \cref{cert:record:analytic-stage}, checked
there in exact rational arithmetic; this is the form of the kernel that
\cref{lem:record:density-can} uses.
\end{corollary}

\begin{proof}
With $\beta=\kappa=1/20$ and $\eta=17/100$, so that $\beta/2=1/40$ and
$\beta/6=1/120$, the seven slacks are, over the common denominator $600$,
\begin{gather*}
 \tfrac{51-30-15-5}{600}=\tfrac1{600},\qquad
 \tfrac{102-30-30-5}{600}=\tfrac{37}{600},\qquad
 \tfrac{600-102-60-45-5}{600}=\tfrac{388}{600}=\tfrac{97}{150},\\
 \tfrac{300-204-60-30-5}{600}=\tfrac1{600},\qquad
 1-\tfrac{51}{100}=\tfrac{49}{100},\qquad
 \tfrac{17}{100}-\tfrac3{40}=\tfrac{19}{200},\qquad
 1-\tfrac{34}{100}-\tfrac3{40}=\tfrac{117}{200}.
\end{gather*}
All are positive, and $\beta/6=1/120$.
\end{proof}

\paragraph{How square-difference-freeness enters.}
The multiplier $g(\xi)$ is the average of $e(\xi m^2)$ over the units modulo
$q(\xi)$. Three hypotheses of the proposition are the form in which
square-difference-freeness is used: $F$ and $J$ are supported in single
residue classes $a$ and $b$ modulo $D$ whose difference is a nonvanishing
square class (strict in (i), of controlled valuation in (ii)); every $x\in\supp F$ is
smaller than every $y\in\supp J$; and no difference $y-x$ is the square of a
positive integer. In the applications $F$ and $J$ are built from a
square-difference-free set $A$, restricted to two adjacent intervals and to
two residue classes that differ by a square, and this makes the third
hypothesis true. The proof compares two kernels on the positive integers: a
model kernel $k_0$, the multiplier $\sum_{\xi\in\Xi_{H,D}}g(\xi)e(-\xi n)$
restricted to $1\le n\le N$ and to the class $n\equiv b-a\pmod D$, and an
actual-square kernel $k$, supported on the squares $m^2$, $1\le m\le\sqrt N$,
weighted by the root-class weights of period dividing $D$ and by the
truncated indicator, expanded in additive characters in
\cref{def:record:K20-kernels}, that the root is a unit at every prime
outside $D$ up to the cutoff $H$. The root classes of \cref{lem:kernel-depth-classes}
give one construction for all three parts of the proposition. Put
$B_h(F,J)=N^{-1}\sum_{x,n}h(n)F(x)J(x+n)$. The ordering and the class condition
give $B_{k_0}(F,J)=D\sum_{\xi\in\Xi_{H,D}}g(\xi)\widehat F(-\xi)\widehat J(\xi)$,
the pairing of the proposition, while $B_k(F,J)=0$ because no difference of
the supports is a positive square. Plancherel bounds $|B_{k_0-k}(F,J)|$ by
$N^{-1}\sup_\alpha|S_0(\alpha)-S(\alpha)|\,\|F\|_{\ell^2}\|J\|_{\ell^2}$, where
$S_0(\alpha)=\sum_nk_0(n)e(\alpha n)$ and
$S(\alpha)=\sum_nk(n)e(\alpha n)$ are the Fourier series of the
two kernels, and the circle method shows that this supremum is
$O(H^{-1/3})$ with a constant depending only on
$\beta,\kappa,\eta,C_D,C_H$. This is the only place in the
analytic stage where the defining property of $A$ yields a gain. The argument is a
quadratic circle-method input in the tradition of Davenport and Vaughan
\cite{Davenport00,Vaughan97}; for the large-sieve and Fourier background see
also \cite{MontgomeryVaughan,IwaniecKowalski}. It follows
\cite[Prop.~6.1]{OpenAI}, with prime-power moduli, a modulus and a cutoff that
are only comparable to powers of $N$, the division of (ii) and symbolic
exponents.

\paragraph{The $\ell^2$ form and the division.}
No bound on $|F|$ or $|J|$ is needed for \eqref{eq:record:K20}. The $\ell^2$
norms make the estimate homogeneous of degree one in each input, which is what
lets the pair cancellation of \cref{lem:record:density-can} retain the
counts of $A$ in the two residue classes in the two intervals, and so gives
a remainder bounded by a constant multiple of the density. A zero
difference is not covered by (ii), and it does not occur under
\ref{it:record:R5}: the exceptional clique records, on which a query difference
vanishes, are removed from the comparison law, and \cref{lem:frontier:iidenergy}
bounds the price of this removal, which is counted in the symmetric part only. Every divided prime keeps an exponent
$2m_p-2j_p\ge2$ in $D_R$, so $D_R$ has the same prime divisors as $D$ and the
frequency set is unchanged, $\Xi_{H,D_R}=\Xi_{H,D}$. The factor $R(N'/N)^2$ in
(ii) is essential: without it the two pairings differ by the factor
$R(N'/N)^2$ of (ii) itself, and equality of the ratios $D/|I|$ and
$(D/R)/(|I|/R)$ for an interval $I$ and its rescaling alone is not the
normalized Fourier identity of (ii).

\paragraph{The range of the exponents.}
With $-\beta/6$ moved to the left, each of the seven conditions says that an
affine function of $(\beta,\kappa)$ with nonnegative coefficients is smaller
than a number depending only on $\eta$. So with $\kappa=\beta$ and
$\eta=17/100$ they hold for every $\beta\in(0,1/20]$, since they hold at
$\beta=1/20$: the proposition contains the kernel at every smaller cutoff
exponent, in particular at $\beta\le1/36$ with squarefree $d$. For
$\kappa=\beta$ the admissible $\eta$ form the window
$(10\beta/3,\,1/4-19\beta/12)$, nonempty only for $\beta<3/59$
(\cref{rem:record:K20-range}); so $\beta=1/20$ is close to the largest cutoff
exponent that this method allows when the modulus is as large as the cutoff.

\paragraph{Consequences and what to read next.}
This section has used the strict square classes of
\cref{def:strict-square-class} and the product formula of
\cref{lem:square-shift}(d) only in its definitions. The proof of
the proposition is \cref{app:kernel20}.
